% open-logic.tex
% open-logic-dev biçemiyle eksiksiz metni üretir

\documentclass[../include/open-logic-part]{subfiles}

\begin{document}

\clearpage

\begin{editorial}
Bu dosya, Open Logic Project kapsamındaki bütün içeriği yükler. Buradaki gibi
yayın kurulu notlarının görüntülenmesi, dosyanın bu malzemenin nasıl sunulacağı
düşünülmeden derlendiğini gösterir. Bunu okuyabiliyorsanız söz konusu PDF'den
ders vermek ya da çalışmak büyük olasılıkla \emph{önerilmez}.

Open Logic Project, öğretime veya kendi kendine çalışmaya daha uygun bir metin üretmeye
yarayan birçok düzenek sunar. Örneğin metin, varsayılan ayarlarda bütün
mantıksal işleçleri ilkel kabul eder ve ispatlarda bütün işleçlerin bütün
durumlarını ele alır. Oysa bu durumlardan bazılarının alıştırma olarak
bırakılması çok daha iyidir. Open Logic Project ayrıca sürmekte olan bir
çalışmadır. İş birliğini ve gelişmeyi özendirmek amacıyla, henüz yalnızca taslak
hâlinde olan veya alıştırmaları eksik bulunan malzeme de projeye katılır. Bir
ders için üretilen PDF bu kesimleri dışarıda bırakacaktır.

Öğretim ve çalışma için daha elverişli PDF'ler bulmak üzere OLP web sitesinde
sunulan
\href{http://builds.openlogicproject.org/}{örnek derslere bakabilirsiniz}.
Kendi dersinizi hazırlamak için
\href{https://github.com/OpenLogicProject/OpenLogic/tree/master/courses/sample}{örnek
  sürücü dosyasından} başlayabilir veya daha ayrıntılı ve ileri örnekler için
türetilmiş ders kitaplarının kaynaklarını inceleyebilirsiniz.
\end{editorial}

\olimport[sets-functions-relations]{sets-functions-relations-complete}

\olimport[propositional-logic]{propositional-logic}

\olimport[first-order-logic]{first-order-logic}

\olimport[model-theory]{model-theory}

\olimport[computability]{computability}

\olimport[turing-machines]{turing-machines}

\olimport[incompleteness]{incompleteness}

\olimport[second-order-logic]{second-order-logic}

\olimport[lambda-calculus]{lambda-calculus}

\olimport[many-valued-logic]{many-valued-logic}

\olimport[normal-modal-logic]{normal-modal-logic}

\olimport[applied-modal-logic]{applied-modal-logic}

\olimport[intuitionistic-logic]{intuitionistic-logic}

\olimport[counterfactuals]{counterfactuals}

\olimport[set-theory]{set-theory}

\olimport[methods]{methods}

\olimport[history]{history}

\olimport[reference]{reference}

\end{document}
